\documentclass[aspectratio=1610]{beamer}

% Prezentacija pripremljena za wildcat beamer temu.
% Kompajliranje: xelatex sertifikati_prezentacija_sa_vise_teksta.tex

\usepackage{fontspec}
\usepackage[serbian]{babel}
\usepackage{amsmath,amssymb,mathtools}
\usepackage{booktabs}
\usepackage{tikz}
\usetikzlibrary{arrows.meta,positioning,calc,shapes.misc}

\title{Digitalni sertifikati i protokoli za uspostavljanje ključa}
\author{Ana Knežević}
\date{\today}

\usetheme[sectionpage=none]{wildcat}

% -----------------------------------------------------------------------------
% Podešavanja izgleda
% -----------------------------------------------------------------------------
\setbeamertemplate{navigation symbols}{}
\setbeamertemplate{itemize items}[circle]
\setbeamertemplate{enumerate items}[default]
\setlength{\abovedisplayskip}{0.30em}
\setlength{\belowdisplayskip}{0.30em}
\setbeamersize{text margin left=0.45cm,text margin right=0.45cm}

\newcommand{\enc}[2]{\{#1\}_{#2}}
\newcommand{\sig}[2]{\mathsf{Sig}_{#1}(#2)}
\newcommand{\Enc}[2]{\mathsf{Enc}_{#1}(#2)}
\newcommand{\Hash}[1]{H(#1)}
\newcommand{\pk}{\mathsf{pk}}
\newcommand{\sk}{\mathsf{sk}}
\newcommand{\Cert}{\mathsf{Cert}}
\newcommand{\CRL}{\mathsf{CRL}}
\newcommand{\serial}{\mathsf{serial}}

\tikzset{
    entity/.style={draw, rounded corners, thick, minimum width=1.35cm, minimum height=0.75cm, align=center},
    bigentity/.style={draw, rounded corners, thick, minimum width=2.7cm, minimum height=0.8cm, align=center},
    attacker/.style={draw, circle, thick, minimum size=0.75cm},
    msg/.style={-Latex, thick},
    trust/.style={-Latex, thick, dashed}
}

\begin{document}

% -----------------------------------------------------------------------------
\begin{frame}
    \titlepage
\end{frame}

\begin{frame}[t]{Pregled}
    \tableofcontents
\end{frame}

% -----------------------------------------------------------------------------
\section{Uvod}

\begin{frame}[t]{Problem uspostavljanja ključa}
\small

\begin{block}{Polazni problem}
U simetričnoj kriptografiji obe strane moraju da imaju isti tajni ključ. 
Ako strane komuniciraju preko nesigurne mreže, postavlja se pitanje kako da bezbedno uspostave taj ključ.
\[
A, B \leadsto K_{AB}
\]
\end{block}

\vspace{0.35em}

\begin{columns}[T,onlytextwidth]
\column{0.50\textwidth}
\begin{alertblock}{Glavni rizik}
Javni ključ je dostupan svima. Napadač može objaviti svoj javni ključ kao da pripada strani $B$.
\end{alertblock}

\column{0.46\textwidth}
\centering
\begin{tikzpicture}[node distance=1.9cm,scale=0.9,transform shape]
    \node[entity] (A) {$A$};
    \node[entity,right=of A] (B) {$B$};
    \node[attacker,below=0.75cm of $(A)!0.5!(B)$] (E) {$E$};
    \draw[msg] (A) -- node[above] {poruke} (B);
    \draw[trust] (E) -- node[right] {prisluškuje} ($(A)!0.5!(B)$);
\end{tikzpicture}
\end{columns}

\vspace{0.45em}

\begin{block}{Zato je potreban sertifikat}
Javni ključ mora biti pouzdano povezan sa identitetom subjekta. 
Nije dovoljno dobiti tajni ključ; učesnik mora znati \textbf{sa kim} ga deli.
\end{block}

\end{frame}

\begin{frame}[t]{Dugoročni i kratkoročni ključevi}
\small
\begin{columns}[T,onlytextwidth]
\column{0.48\textwidth}
\begin{block}{Dugoročni ključevi}
\begin{itemize}
    \item koriste se duže vreme;
    \item povezani su sa identitetom učesnika;
    \item primer: privatni ključ subjekta čiji se javni ključ nalazi u sertifikatu;
    \item kompromitovanje narušava poverenje u identitet.
\end{itemize}
\end{block}

\column{0.48\textwidth}
\begin{block}{Kratkoročni ključevi}
\begin{itemize}
    \item sesijski ili efemerni ključevi;
    \item koriste se za jednu sesiju ili ograničeno vreme;
    \item ograničavaju posledice napada;
    \item cilj je da se bezbedno uspostave pomoću dugoročnih ključeva.
\end{itemize}
\end{block}
\end{columns}

\vspace{0.5em}
\begin{alertblock}{Bezbednosni zahtevi}
\centering
\textbf{tajnost} \quad + \quad \textbf{autentifikacija} \quad + \quad \textbf{svežina} \quad + \quad \textbf{ispravno vezivanje}
\end{alertblock}

\vspace{0.2em}
\centering
\footnotesize
Ključ treba da bude tajan, nov, vezan za prave učesnike i za konkretnu sesiju.
\end{frame}

% -----------------------------------------------------------------------------
\section{Sertifikati i sertifikaciona tela}

\begin{frame}[t]{Digitalni sertifikat}
\small
\begin{columns}[T,onlytextwidth]
\column{0.52\textwidth}
\begin{block}{Definicija}
Digitalni sertifikat je potpisana struktura koja povezuje javni ključ sa identitetom subjekta.
\end{block}

\begin{itemize}
    \item subjekt može biti osoba, server, organizacija ili proces;
    \item CA proverava identitet i potvrđuje vezu sa javnim ključem;
    \item korisnik proverava potpis sertifikacionog tela.
\end{itemize}

\begin{block}{Ideja}
\[
\Cert_A \;:\; A \longleftrightarrow \pk_A
\]
\centering
sertifikat kaže: „ovaj javni ključ pripada subjektu $A$”
\end{block}

\column{0.44\textwidth}
\centering
\begin{tikzpicture}
    \node[draw, rounded corners, thick, minimum width=4.25cm, minimum height=3.1cm, align=left] (cert) {};
    \node[anchor=north west] at ([xshift=0.18cm,yshift=-0.15cm]cert.north west) {\textbf{Sertifikat}};
    \node[anchor=west] at ([xshift=0.25cm,yshift=0.45cm]cert.west) {Subject: $A$};
    \node[anchor=west] at ([xshift=0.25cm,yshift=0.00cm]cert.west) {Public key: $\pk_A$};
    \node[anchor=west] at ([xshift=0.25cm,yshift=-0.45cm]cert.west) {Issuer: $CA$};
    \node[anchor=west] at ([xshift=0.25cm,yshift=-0.90cm]cert.west) {Signature: $\sig{\sk_{CA}}{\cdots}$};
\end{tikzpicture}

\vspace{0.6em}
\begin{alertblock}{Zašto je CA potreban?}
Bez CA svako bi mogao da tvrdi da neki javni ključ pripada nekom subjektu.
\end{alertblock}
\end{columns}
\end{frame}

\begin{frame}[t]{Provera sertifikata}
\small
\begin{columns}[T,onlytextwidth]
\column{0.50\textwidth}
\begin{block}{Korisnik proverava}
\begin{itemize}
    \item potpis izdavaoca;
    \item period važenja;
    \item namenu javnog ključa;
    \item ograničenja sertifikata;
    \item status opoziva.
\end{itemize}
\end{block}

\column{0.46\textwidth}
\begin{block}{Provera u praksi}
\begin{itemize}
    \item da li je sertifikat zaista izdao navedeni CA;
    \item da li sertifikat trenutno važi;
    \item da li se ključ sme koristiti za tu namenu;
    \item da li je sertifikat opozvan.
\end{itemize}
\end{block}
\end{columns}

\vspace{0.3em}
\begin{alertblock}{Važno}
Sertifikat potvrđuje javni ključ, ali ne štiti privatni ključ subjekta. Ako privatni ključ procuri, sertifikat može biti zloupotrebljen.
\end{alertblock}
\end{frame}

\begin{frame}[t]{Lanac poverenja}
\small
\centering
\begin{tikzpicture}[node distance=1.45cm]
    \node[bigentity] (root) {Root CA};
    \node[bigentity,right=of root] (inter) {Intermediate CA};
    \node[bigentity,right=of inter] (end) {Krajnji subjekt};
    \draw[msg] (root) -- node[above] {potpisuje} (inter);
    \draw[msg] (inter) -- node[above] {potpisuje} (end);
\end{tikzpicture}

\vspace{0.6em}
\begin{columns}[T,onlytextwidth]
\column{0.49\textwidth}
\begin{block}{Certification path}
\[
CA_0 \rightarrow CA_1 \rightarrow \cdots \rightarrow A
\]
Poverenje se prenosi od poznatog CA do krajnjeg subjekta.
\end{block}

\column{0.47\textwidth}
\begin{block}{Sidro poverenja}
Korisnik unapred veruje javnom ključu korenskog CA. Od tog javnog ključa počinje provera lanca.
\end{block}
\end{columns}

\vspace{0.25em}
\begin{alertblock}{Validnost lanca}
Svi potpisi, periodi važenja, opozivi i ograničenja moraju biti ispravni. U većim sistemima poverenje se može proširiti i unakrsnom sertifikacijom.
\end{alertblock}
\end{frame}

\begin{frame}[t]{X.509 i opoziv sertifikata}
\small
\begin{columns}[T,onlytextwidth]
\column{0.49\textwidth}
\begin{block}{Tipična polja X.509 sertifikata}
\begin{itemize}
    \item Subject: kome sertifikat pripada;
    \item Issuer: ko ga je izdao;
    \item javni ključ subjekta;
    \item serijski broj;
    \item period važenja;
    \item digitalni potpis;
    \item v3 ekstenzije.
\end{itemize}
\end{block}

\column{0.47\textwidth}
\begin{block}{CRL}
CRL je potpisana lista opozvanih sertifikata.

\medskip
Sertifikat može biti opozvan i pre isteka važenja.

\medskip
Svaki opozvani sertifikat se u listi prepoznaje po svom serijskom broju.
\end{block}

\begin{alertblock}{Razlog opoziva}
Kompromitovan ključ, promenjeni podaci ili narušeno poverenje.
\end{alertblock}
\end{columns}
\end{frame}

% -----------------------------------------------------------------------------
\section{Sesijski ključevi iz statičkih simetričnih ključeva}

\begin{frame}[t]{Model sa poverljivom trećom stranom}
\small
\begin{columns}[T,onlytextwidth]
\column{0.50\textwidth}
\begin{block}{Pretpostavke}
\begin{itemize}
    \item $S$ je server za distribuciju ključeva;
    \item $A$ i $S$ dele dugoročni ključ $K_{AS}$;
    \item $B$ i $S$ dele dugoročni ključ $K_{BS}$;
    \item cilj je novi sesijski ključ $K_{AB}$.
\end{itemize}
\end{block}

\begin{block}{Oznaka}
\[
\enc{M}{K}
\]
\centering
poruka $M$ šifrovana ključem $K$
\end{block}

\column{0.46\textwidth}
\centering
\begin{tikzpicture}[node distance=1.55cm]
    \node[entity] (A) {$A$};
    \node[entity,right=2.7cm of A] (B) {$B$};
    \node[entity,below=1.15cm of $(A)!0.5!(B)$] (S) {$S$};
    \draw[msg,<->] (A) -- node[left] {$K_{AS}$} (S);
    \draw[msg,<->] (B) -- node[right] {$K_{BS}$} (S);
    \draw[trust,<->] (A) -- node[above] {$K_{AB}$} (B);
\end{tikzpicture}

\vspace{0.7em}
\begin{alertblock}{Uloga servera}
$S$ pomaže da $A$ i $B$ dobiju novi ključ, ali ne učestvuje u kasnijoj komunikaciji.
\end{alertblock}
\end{columns}
\end{frame}

\begin{frame}[t]{Wide-Mouth Frog protokol}
\small
\begin{block}{Ideja}
Strana $A$ sama bira sesijski ključ $K_{AB}$ i preko servera $S$ ga dostavlja strani $B$.
\end{block}

\begin{block}{Tok poruka}
\[
\begin{aligned}
1.\quad A &\rightarrow S : A,\ \enc{T_A, B, K_{AB}}{K_{AS}} \\
2.\quad S &\rightarrow B : \enc{T_S, A, K_{AB}}{K_{BS}}
\end{aligned}
\]
\end{block}

\begin{columns}[T,onlytextwidth]
\column{0.48\textwidth}
\begin{block}{Značenje poruka}
\begin{itemize}
    \item $T_A$ i $T_S$ služe za proveru svežine;
    \item $B$ dobija ključ preko servera $S$;
    \item server samo prosleđuje ključ koji bira $A$.
\end{itemize}
\end{block}

\column{0.48\textwidth}
\begin{alertblock}{Ograničenje}
Protokol zavisi od sinhronizacije satova i od kvaliteta ključa koji bira $A$.
\end{alertblock}
\end{columns}
\end{frame}

\begin{frame}[t]{Needham--Schroeder protokol}
\scriptsize
\begin{block}{Ideja}
Sesijski ključ ne bira $A$, već ga generiše server $S$. Nonce $N_A$ povezuje zahtev i odgovor servera.
\end{block}

\begin{block}{Tok poruka}
\[
\begin{aligned}
1.\quad A &\rightarrow S : A, B, N_A \\
2.\quad S &\rightarrow A : \enc{N_A, B, K_{AB}, \enc{K_{AB}, A}{K_{BS}}}{K_{AS}} \\
3.\quad A &\rightarrow B : \enc{K_{AB}, A}{K_{BS}} \\
4.\quad B &\rightarrow A : \enc{N_B}{K_{AB}} \\
5.\quad A &\rightarrow B : \enc{N_B - 1}{K_{AB}}
\end{aligned}
\]
\end{block}

\end{frame}

\begin{frame}[t]{Needham--Schroeder: značenje poruka}
\small
\begin{columns}[T,onlytextwidth]
\column{0.49\textwidth}
\begin{block}{Od servera do učesnika}
\begin{itemize}
    \item $A$ šalje $N_A$ da bi kasnije prepoznala odgovor.
    \item Server $S$ generiše novi $K_{AB}$.
    \item Deo šifrovan ključem $K_{BS}$ namenjen je strani $B$.
\end{itemize}
\end{block}

\column{0.47\textwidth}
\begin{block}{Potvrda posedovanja ključa}
\begin{itemize}
    \item $B$ šalje $N_B$ šifrovan ključem $K_{AB}$.
    \item $A$ vraća $N_B-1$.
    \item Time $A$ pokazuje da poseduje isti sesijski ključ.
\end{itemize}
\end{block}
\end{columns}

\vspace{0.3em}
\begin{alertblock}{Slabost}
$B$ nema direktan dokaz da je $K_{AB}$ nov, pa je svežina ključna slabost protokola. Ako se kasnije otkrije stari sesijski ključ, napadač može ponoviti staru poruku namenjenu strani $B$.
\end{alertblock}
\end{frame}

\begin{frame}[t]{Kerberos}
\scriptsize
\begin{block}{Ideja}
Kerberos koristi tikete i vremenske oznake. Tiket je deo poruke namenjen strani $B$, koji $A$ samo prosleđuje.
\end{block}

\begin{block}{Tok poruka}
\[
\begin{aligned}
1.\quad A &\rightarrow S : A, B \\
2.\quad S &\rightarrow A : \enc{T_S, L, K_{AB}, B, \enc{T_S, L, K_{AB}, A}{K_{BS}}}{K_{AS}} \\
3.\quad A &\rightarrow B : \enc{T_S, L, K_{AB}, A}{K_{BS}},\ \enc{A,T_A}{K_{AB}} \\
4.\quad B &\rightarrow A : \enc{T_A+1}{K_{AB}}
\end{aligned}
\]
\end{block}


\end{frame}

\begin{frame}[t]{Kerberos: tiket i potvrda}
\small
\begin{columns}[T,onlytextwidth]
\column{0.48\textwidth}
\begin{block}{Tiket za $B$}
\[
\enc{T_S,L,K_{AB},A}{K_{BS}}
\]
$A$ ne može da menja tiket, već ga samo pokazuje strani $B$.
\end{block}

\column{0.48\textwidth}
\begin{block}{Potvrda učesnika}
\[
\enc{A,T_A}{K_{AB}}
\]
Ova poruka pokazuje da $A$ zna sesijski ključ i da trenutno učestvuje u komunikaciji.
\end{block}
\end{columns}

\begin{alertblock}{Prednost i ograničenje}
Koristi tikete sa ograničenim vremenom važenja $L$. Time se smanjuje mogućnost ponavljanja starih poruka. 
Zahteva sinhronizaciju kljuceva. Ako vremenska oznaka nije sveža, poruka se odbacuje.
\end{alertblock}

\end{frame}

% -----------------------------------------------------------------------------
\section{Sesijski ključevi iz statičkih javnih ključeva}

\begin{frame}[t]{Prenos ključa pomoću javnog ključa}
\small
\begin{columns}[T,onlytextwidth]
\column{0.50\textwidth}
\begin{block}{Oznake}
\begin{itemize}
    \item $\pk_B$ je javni ključ strane $B$;
    \item $\sk_B$ je privatni ključ strane $B$;
    \item $K_{AB}$ je sesijski ključ koji bira $A$.
\end{itemize}
\end{block}

\begin{block}{Osnovna ideja}
\[
A \rightarrow B : \Enc{\pk_B}{K_{AB}}
\]
Samo $B$ može da dešifruje poruku, jer samo $B$ ima privatni ključ $\sk_B$.
\end{block}

\column{0.46\textwidth}
\begin{alertblock}{Problem}
$B$ može da dešifruje poruku, ali ne zna ko ju je poslao. Napadač takođe može poslati ključ šifrovan javnim ključem $\pk_B$.
\end{alertblock}

\begin{block}{Zaključak}
Prenos ključa mora dodatno da obezbedi autentifikaciju pošiljaoca.
\end{block}
\end{columns}
\end{frame}

\begin{frame}[t]{Prenos ključa: potpis i identitet}
\small
\begin{columns}[T,onlytextwidth]
\column{0.48\textwidth}
\begin{block}{Sa potpisom i identitetom}
\[
c = \Enc{\pk_B}{A\parallel K_{AB}}
\]
Identitet $A$ se uključuje u šifrovani deo da bi ključ bio vezan za pošiljaoca.
\[
A \rightarrow B : c,\ \sig{\sk_A}{c}
\]
Potpis pokazuje da je poruku potpisala strana koja poseduje privatni ključ $\sk_A$.
\end{block}


\end{columns}

\vspace{0.3em}
\begin{alertblock}{Ograničenje}
Ključ bira samo jedna strana i nema forward secrecy.
\end{alertblock}
\end{frame}

\begin{frame}[t]{Forward secrecy}
\small
\begin{block}{Definicija}
Forward secrecy znači da kompromitovanje dugoročnog privatnog ključa ne otkriva ranije uspostavljene sesijske ključeve.
\end{block}

\begin{columns}[T,onlytextwidth]
\column{0.48\textwidth}
\begin{alertblock}{Prenos ključa}
\centering
\[
\begin{array}{c}
\sk_B\ \text{kompromitovan} \\
\Downarrow \\
K_{AB}\ \text{može biti otkriven}
\end{array}
\]
\smallskip
\small Snimljeni šifrat može se kasnije dešifrovati.
\end{alertblock}

\column{0.48\textwidth}
\begin{block}{Efemerni Diffie--Hellman}
\centering
\[
\begin{array}{c}
a,b\ \text{obrisani} \\
\Downarrow \\
g^{ab}\ \text{ostaje zaštićen}
\end{array}
\]
\smallskip
\small Prošle sesije ostaju zaštićene ako su efemerne tajne obrisane.
\end{block}
\end{columns}
\end{frame}

\begin{frame}[t]{Diffie--Hellman protokol}
\small
\begin{columns}[T,onlytextwidth]
\column{0.50\textwidth}
Diffie--Hellman je osnovni primer dogovora ključa.
\begin{block}{1. Javne vrednosti i tajne}
\begin{itemize}
    \item $G$ je grupa reda $q$;
    \item $g$ je generator;
    \item $A$ bira tajnu $a$;
    \item $B$ bira tajnu $b$.
\end{itemize}
\end{block}

\begin{block}{2. Razmena}
\[
A \rightarrow B : g^a
\]
\[
B \rightarrow A : g^b
\]
\end{block}

\column{0.46\textwidth}
\begin{block}{3. Računanje zajedničke tajne}
\[
A: (g^b)^a=g^{ab}
\]
\[
B: (g^a)^b=g^{ab}
\]
\[
K_{AB}=\Hash{g^{ab}}
\]
\end{block}

\begin{alertblock}{Napomena}
Osnovni DH ne obezbeđuje autentifikaciju. Strane dobijaju ključ, ali ne znaju pouzdano sa kim su ga dobile.
\end{alertblock}
\end{columns}
\end{frame}

\begin{frame}[t]{Napad čoveka u sredini}
\small
\begin{columns}[T,onlytextwidth]
\column{0.52\textwidth}
\begin{block}{Zamena efemernih vrednosti}
\[
\begin{aligned}
1.\quad A &\rightarrow E : g^a \\
2.\quad E &\rightarrow A : g^m \\
3.\quad E &\rightarrow B : g^n \\
4.\quad B &\rightarrow E : g^b
\end{aligned}
\]
\end{block}

\begin{itemize}
    \item $A$ misli da dobija vrednost od $B$;
    \item $B$ misli da dobija vrednost od $A$;
    \item zapravo oba učesnika računaju ključ sa napadačem.
\end{itemize}

\column{0.44\textwidth}
\begin{alertblock}{Posledica}
\[
K_{AE}=\Hash{g^{am}}
\]
\[
K_{BE}=\Hash{g^{bn}}
\]
Osnovni Diffie--Hellman rešava problem tajnog dogovora ključa, ali ne rešava problem autentifikacije.
\end{alertblock}
\end{columns}
\end{frame}

\begin{frame}[t]{Potpisani Diffie--Hellman}
\small
\begin{block}{Oznake}
$a$ i $b$ su efemerne tajne, a $\sk_A$ i $\sk_B$ su dugoročni privatni ključevi za potpis.
\end{block}

\begin{block}{Poruke}
\[
\begin{aligned}
A &\rightarrow B : g^a,\ \sig{\sk_A}{g^a} \\
B &\rightarrow A : g^b,\ \sig{\sk_B}{g^b}
\end{aligned}
\]
\end{block}

\begin{alertblock}{Suština}
Potpis autentifikuje poslatu Diffie--Hellman vrednost. Još bolje je da potpis obuhvati ceo kontekst sesije: obe efemerne vrednosti i uloge učesnika.
\end{alertblock}
\end{frame}

\begin{frame}[t]{STS protokol}
\scriptsize
\begin{block}{Oznake}
$et_A=g^a$ i $et_B=g^b$ su efemerne javne vrednosti. Ključ $k$ se izvodi iz zajedničke Diffie--Hellman tajne.
\end{block}

\begin{block}{Tok poruka}
\[
\begin{aligned}
1.\quad A &\rightarrow B : et_A=g^a \\
2.\quad B &\rightarrow A : et_B=g^b,\ \enc{\sig{\sk_B}{et_B,\,et_A}}{k} \\
3.\quad A &\rightarrow B : \enc{\sig{\sk_A}{et_A,\,et_B}}{k}
\end{aligned}
\]
\[
k=\Hash{g^{ab}}
\]
\end{block}

\begin{alertblock}{Suština}
Potpisi obuhvataju obe efemerne vrednosti, pa napadač ne može jednostavno da ih zameni i zadrži ispravan potpis.
\end{alertblock}
\end{frame}

\begin{frame}[t]{Blake--Wilson--Menezes protokol}
\small
\begin{block}{Ideja}
Tok poruka je kao kod Diffie--Hellmana, ali se ključ izvodi i iz statičkih i iz efemernih vrednosti.
\end{block}

\begin{columns}[T,onlytextwidth]
\column{0.48\textwidth}
\begin{block}{Statički ključevi}
\[
\pk_A=g^a,\qquad \pk_B=g^b
\]
$a$ i $b$ su statički privatni ključevi.
\end{block}

\column{0.48\textwidth}
\begin{block}{Efemerni ključevi}
$x$ i $y$ su nove tajne za konkretnu sesiju.
\[
A\rightarrow B: ek_A=g^x
\]
\[
B\rightarrow A: ek_B=g^y
\]
\end{block}
\end{columns}

\vspace{0.3em}
\begin{alertblock}{Suština}
Statički ključevi vezuju rezultat za identitete, a efemerni ključevi za konkretnu sesiju.
\end{alertblock}
\end{frame}

\begin{frame}[t]{Blake--Wilson--Menezes: računanje ključa}
\small
\begin{columns}[T,onlytextwidth]
\column{0.48\textwidth}
\begin{block}{Računa $A$}
\[
k=H(\pk_B^x,\ ek_B^a)
\]
\end{block}

\begin{block}{Računa $B$}
\[
k=H(ek_A^b,\ \pk_A^y)
\]
\end{block}

\column{0.48\textwidth}
\begin{block}{Zašto je isti ključ?}
\[
\pk_B^x=(g^b)^x=g^{bx}=ek_A^b
\]
\[
ek_B^a=(g^y)^a=g^{ay}=\pk_A^y
\]
\end{block}
\end{columns}

\vspace{0.3em}
\begin{alertblock}{Zaključak}
Obe strane dobijaju iste ulaze za funkciju $H$, pa dobijaju isti ključ.
\end{alertblock}
\end{frame}

\begin{frame}[t]{MQV protokol}
\small
\begin{block}{Ideja}
MQV kombinuje efemerne i statičke ključeve bez dodatnih potpisa i bez treće poruke.
\end{block}

\begin{columns}[T,onlytextwidth]
\column{0.48\textwidth}
\begin{block}{Statički ključevi}
\[
\pk_A=g^a,\qquad \pk_B=g^b
\]
$a$ i $b$ su statički privatni ključevi.
\end{block}

\column{0.48\textwidth}
\begin{block}{Efemerni ključevi}
$x$ i $y$ su efemerne tajne, a $ek_A$ i $ek_B$ efemerne javne vrednosti.
\[
A\rightarrow B: ek_A=g^x,\qquad B\rightarrow A: ek_B=g^y
\]
\end{block}
\end{columns}

\begin{alertblock}{Pomoćne vrednosti}
$s_A,t_A,s_B,t_B$ izvode se iz efemernih javnih ključeva i koriste se pri računanju konačne zajedničke vrednosti.
\end{alertblock}
\end{frame}

\begin{frame}[t]{MQV: računanje ključa}
\small
\begin{columns}[T,onlytextwidth]
\column{0.48\textwidth}
\begin{block}{Strana $A$}
\[
\begin{aligned}
h_A &= x+s_Aa \pmod q,\\
K_A &= (ek_B\cdot \pk_B^{t_A})^{h_A}
\end{aligned}
\]
\end{block}

\column{0.48\textwidth}
\begin{block}{Strana $B$}
\[
\begin{aligned}
h_B &= y+s_Bb \pmod q,\\
K_B &= (ek_A\cdot \pk_A^{t_B})^{h_B}
\end{aligned}
\]
\end{block}
\end{columns}

\vspace{0.3em}
\begin{alertblock}{Suština}
Efemerni ključevi daju svežinu, a statički ključevi vezuju rezultat za identitete učesnika.
\end{alertblock}
\end{frame}

% -----------------------------------------------------------------------------

\section{Simbolička analiza protokola}

\begin{frame}[t]{Zašto simbolička analiza?}
\small

\begin{block}{Polazna ideja}
Protokol može koristiti jake algoritme, a ipak biti nebezbedan zbog pogrešnog toka poruka.
\end{block}

\vspace{0.25em}

\begin{block}{Gde nastaje problem?}
\begin{itemize}
    \item poruka može biti ispravno šifrovana, ali stara;
    \item poruka može biti validna, ali namenjena drugoj sesiji;
    \item učesnik može pogrešno zaključiti šta poruka dokazuje.
\end{itemize}
\end{block}

\vspace{0.25em}

\begin{alertblock}{Zaključak}
Zato nije dovoljno proveriti samo da li je poruka pravilno šifrovana. Važno je proveriti šta ta poruka zaista znači u toku protokola.
\end{alertblock}

\end{frame}

\begin{frame}[t]{Model napadača i cilj analize}
\small

\begin{block}{Model napadača}
Napadač se posmatra kao jak protivnik koji kontroliše mrežu. On može da presreće, menja, ponavlja, zaustavlja i šalje poruke.
\end{block}

\vspace{0.25em}

\begin{block}{Šta se analizira?}
Analizira se logika protokola:
\begin{itemize}
    \item šta učesnik vidi;
    \item u šta učesnik veruje;
    \item šta sme da zaključi nakon prijema poruke.
\end{itemize}
\end{block}

\vspace{0.25em}

\begin{alertblock}{Cilj}
Kod protokola za uspostavljanje ključa cilj je da učesnici na kraju veruju da dele isti svež sesijski ključ.
\end{alertblock}

\end{frame}





\begin{frame}[t]{BAN logika: osnovne oznake}
\small
\begin{columns}[T,onlytextwidth]
\column{0.48\textwidth}
Formalni pristup za analizu autentifikacionih protokola.
\begin{block}{Verovanja i poruke}
\[
P\mid\equiv X \quad \text{$P$ veruje u $X$}
\]
\[
P\triangleleft X \quad \text{$P$ vidi $X$}
\]
\[
P\mid\sim X \quad \text{$P$ je nekada rekao $X$}
\]
\end{block}

\column{0.48\textwidth}
\begin{block}{Svežina i ključevi}
\[
\#(X) \quad \text{$X$ je sveže}
\]
\[
P\Rightarrow X \quad \text{$P$ je nadležan za $X$}
\]
\[
P \xleftrightarrow{K} Q \quad \text{$P$ i $Q$ dele tajni ključ $K$}
\]
\end{block}
\end{columns}

\vspace{0.3em}
\begin{alertblock}{Upotreba oznaka}
Oznake služe da formalno opišemo šta učesnik vidi, u šta veruje i šta sme da zaključi nakon prijema poruke.
\end{alertblock}
\end{frame}

\begin{frame}[t]{BAN logika: cilj protokola}
\small
\begin{block}{Cilj protokola za ključ}
\[
A\mid\equiv A\xleftrightarrow{K_{AB}}B,
\qquad
B\mid\equiv A\xleftrightarrow{K_{AB}}B
\]
Obe strane treba da veruju da dele isti sesijski ključ.
\end{block}

\begin{alertblock}{Ideja analize}
Ako se cilj ne može izvesti iz poruka i pretpostavki, protokol ima slabost ili nije dovoljno precizno definisan.
\end{alertblock}
\end{frame}

\begin{frame}[t]{BAN logika: pravila zaključivanja}
\scriptsize
\begin{columns}[T,onlytextwidth]
\column{0.48\textwidth}
\begin{block}{Značenje poruke}
\[
\frac{A\mid\equiv A\xleftrightarrow{K}B,\quad A\triangleleft \enc{X}{K}}
{A\mid\equiv B\mid\sim X}
\]
Ako $A$ veruje da deli ključ $K$ sa $B$, onda iz poruke šifrovane tim ključem zaključuje da je $B$ nekada rekao $X$.
\end{block}

\begin{block}{Svežina}
\[
\frac{A\mid\equiv \#(X),\quad A\mid\equiv B\mid\sim X}
{A\mid\equiv B\mid\equiv X}
\]
Ako je $X$ sveže, poruka nije samo ponovljena iz stare sesije.
\end{block}

\column{0.48\textwidth}
\begin{block}{Nadležnost}
\[
\frac{A\mid\equiv B\Rightarrow X,\quad A\mid\equiv B\mid\equiv X}
{A\mid\equiv X}
\]
Ako je $B$ nadležan za $X$ i veruje u $X$, onda $A$ može prihvatiti $X$.
\end{block}

\begin{alertblock}{Zaključak}
Pravila povezuju poruke, svežinu i poverenje u autoritet učesnika.
\end{alertblock}
\end{columns}
\end{frame}

% -----------------------------------------------------------------------------
\section{Zaključak}

\begin{frame}[t]{Zaključak}
\small
\begin{columns}[T,onlytextwidth]
\column{0.52\textwidth}
\begin{itemize}
    \item Uspostavljanje ključa nije samo izbor algoritma.
    \item Sertifikati povezuju javni ključ sa identitetom.
    \item CA svojim potpisom potvrđuje tu vezu.
    \item Autentifikovani protokoli sprečavaju napad čoveka u sredini.
    \item Formalna analiza otkriva greške u toku poruka i zaključivanju.
\end{itemize}

\column{0.44\textwidth}
\begin{alertblock}{Suština}
Učesnici ne treba samo da dobiju tajni ključ, već da budu sigurni sa kim ga dele, da je ključ svež i da je namenjen baš toj sesiji.
\end{alertblock}
\end{columns}
\end{frame}

\begin{frame}[t]{Literatura}
\small
\begin{itemize}
    \item Nigel P. Smart, \emph{Cryptography Made Simple}, Springer, 2016.
    \item William Stallings, \emph{Cryptography and Network Security}, 2011.
    \item RFC 5280, \emph{Internet X.509 Public Key Infrastructure Certificate and CRL Profile}.
    \item M. Burrows, M. Abadi, R. Needham, \emph{A Logic of Authentication}, 1990.
    \item IBM Documentation, \emph{Digital certificates}.
\end{itemize}
\end{frame}

\end{document}
